\documentclass[11pt]{article}
\usepackage{amsfonts,amsmath,amssymb,bm,booktabs,enumitem,array}
\usepackage[letterpaper,textwidth=6.3in,textheight=8.7in]{geometry}
\usepackage[T1]{fontenc}
\usepackage{xcolor}
\usepackage{listings}
\usepackage[authoryear,round]{natbib}
\usepackage[breaklinks=true,colorlinks,linkcolor=blue,citecolor=blue,urlcolor=blue]{hyperref}
\hypersetup{pdfpagelayout=SinglePage}% opens one page at a time, not as a continuous scroll

\lstset{basicstyle=\ttfamily\small,columns=fullflexible,keepspaces=true,frame=single,framerule=0.3pt,rulecolor=\color{black!40},xleftmargin=0.5em,framexleftmargin=0.5em,breaklines=true,showstringspaces=false,aboveskip=0.6em,belowskip=0.6em}
\newcommand{\code}[1]{\texttt{#1}}
\newcommand{\E}{{\mathbb E}}
\newcommand{\PP}{{\mathbb P}}
\newcommand{\iu}{{\mathrm i}}
\newcommand{\dr}[1]{{\mathrm d}#1}
\newcommand{\Rea}{\operatorname{Re}}
\newcommand{\phih}{\hat\phi}
\newcommand{\tauh}{\hat\tau}
\newcommand{\Fk}{\textsf{DK}}
\newcommand{\Fu}{\textsf{DU}}
\newcommand{\Ft}{\textsf{DT}}
\newcommand{\Sk}{\textsf{SK}}
\newcommand{\Su}{\textsf{SU}}
\newcommand{\St}{\textsf{ST}}
\setlist[description]{leftmargin=1.5em,itemsep=0.35em,style=nextline,font=\normalfont\ttfamily}
\setlength{\parskip}{0.25em}

\title{The Exact Limiting Distribution of the Dickey--Fuller $t$-Ratio\\ and Some Related Statistics\\[0.6ex]\large MATLAB programs: user manual}
\author{Yong Bao\\\normalsize Purdue University \and Simon A. Broda\\\normalsize Lucerne University of Applied Sciences \and Raymond Kan\\\normalsize University of Toronto and\\\normalsize Hong Kong University of Science and Technology}
\date{Version 0.4.0, November 2026}

\begin{document}
\maketitle

\begin{abstract}
\noindent This manual documents the MATLAB programs that accompany \citet*{BaoBrodaKan2026l} (``the paper'').  The programs compute the limiting distribution function, density, quantiles and moments of the least squares $t$-ratio $\tauh_0=(\phih-\phi_0)/\operatorname{se}(\phih)$ in a Gaussian first-order autoregression under the local alternatives $\phi=1+c/n$, for the regression with no deterministic terms, an intercept, or an intercept and a linear trend, for the null value $\phi_0=1+c_0/n$ ($c_0=0$: the Dickey--Fuller tests), and with a zero, a fixed or a stationary initial observation of the order of $\sqrt n$.  Two exact representations are implemented for every case, the line integral representation and the conditional integral representation of the paper (its first part for a region of the plane, its second part the integration-by-parts form), and they check each other; with a fixed start in the cases with an intercept or a trend a third, the contour integral representation (the inversion formula on paths around the singularities of the moment generating function, with the integral over $v$ in closed form and the paths of the inversion deformed), takes the place of the line integral.  Section~\ref{sec:start} is a quick start, Section~\ref{sec:user} describes the user-level routines, Section~\ref{sec:engines} the engines, Section~\ref{sec:accuracy} accuracy and cost, Section~\ref{sec:validation} the validation suite and the reproduction of the paper, and Section~\ref{sec:files} lists the files.
\end{abstract}

\tableofcontents

\section{Quick start}
\label{sec:start}

Add the folder \code{matlab} of the package to the MATLAB path.  The model is
\[
y_t=\alpha+\delta t+\phi y_{t-1}+\sigma e_t,\qquad t=1,\ldots,n,\qquad e_t\sim N(0,1)\text{ i.i.d.},
\]
with $\phi=1+c/n$, no drift, and $y_0=\sigma\sqrt n\,\gamma_0$.  The case code names the start (\textsf{D}eterministic, that is fixed, or \textsf{S}tationary) and the regressors (intercept \textsf{K}nown to be zero; \textsf{U}nknown intercept; intercept and \textsf{T}rend):
\begin{center}
\begin{tabular}{@{}lll@{}}
\toprule
Code & Regression & Start\\
\midrule
\code{'DK'} & $y_t$ on $y_{t-1}$ ($\hat\tau$) & $\gamma_0$ fixed (\code{opts.g0}, default $0$)\\
\code{'DU'} & $y_t$ on $1$, $y_{t-1}$ ($\hat\tau_\mu$) & $\gamma_0$ fixed\\
\code{'DT'} & $y_t$ on $1$, $t$, $y_{t-1}$ ($\hat\tau_\tau$) & $\gamma_0$ fixed\\
\code{'SK'}, \code{'SU'}, \code{'ST'} & the same & stationary, $J(0)\sim N(0,-1/(2c))$, $c\le-10^{-4}$\\
\bottomrule
\end{tabular}
\end{center}
With a stationary start the routines raise an error for $c>-10^{-4}$: closer to zero the variance $-1/(2c)$ of the start is too large for the quadrature.
All routines take \code{(\ldots,cas,c,opts)}: the case, the local-to-unity parameter $c$ (default $0$, the unit root) and an optional structure of options, of which \code{c0} (the local null value, default $0$), \code{g0} (the fixed start, default $0$) and \code{method} (\code{'auto'}, the default, \code{'single'}, \code{'double'} or, with a fixed start in cases \Fu, \Ft, \code{'contour'}) and \code{verified} (\code{false}, the default: one evaluation by one method; \code{true}: the checked computation, Section~\ref{sec:accuracy}) are the ones a user needs.
\begin{lstlisting}
ltratio_icdf([0.01 0.05 0.10], 'DU')        % asymptotic critical values of tau_mu
ltratio_cdf(-2.86157, 'DU', -5)              % limiting power of the 5% test at c = -5
ltratio_cdf(-2.86157, 'SU', -5)              % the same with a stationary start
ltratio_cdf(-2.86157, 'DU', -5, struct('g0', 1.5))   % fixed start y_0 = 1.5 sigma sqrt(n) (contour integral)
ltratio_cdf(0, 'DK', 2, struct('c0', 2))     % t-ratio centered at the true value, c = 2
ltratio_pdf(-1:0.5:1, 'DK')                  % the density of tau at five points
ltratio_mom(4, 'DT')                         % E tau_tau^j, j = 1..4
\end{lstlisting}
The first call returns $-3.43035$, $-2.86157$ and $-2.56679$; the second $0.1174$.  A value of the distribution function or the density takes a fraction of a second, a quantile a few values, and a set of moments about a second.

\section{The user-level routines}
\label{sec:user}

\paragraph{Common arguments.}  \code{cas} is the case code (in either letter case), \code{c} a finite real scalar (negative for a stationary start), and \code{opts} a structure.  The arguments are checked by \code{ltratio\_args}, which rejects unknown codes, a stationary start with $c\ge0$, a fixed start \code{g0} in a stationary case, an unknown \code{method} and \code{'contour'} without a fixed start in cases \Fu, \Ft.  Fields of \code{opts} other than \code{c0}, \code{g0}, \code{method} and \code{verified} are passed to the engine in use (Section~\ref{sec:engines}).  Vector arguments are evaluated elementwise.

\begin{description}
\item[{[F,info] = ltratio\_cdf(x,cas,c,opts)}]
The limit of $\PP[\tauh_0\le x]$.  At \code{x = -Inf}, \code{Inf} and \code{NaN} the result is $0$, $1$ and \code{NaN}.  \code{info.method} is the representation used and \code{info.time} the elapsed time; with \code{'double'} the diagnostics of \code{tr\_limdist} are added (\code{info.status}, \code{info.err}, \ldots; a nonzero status raises the warning \code{tr\_limdist:status}), with \code{'contour'} those of \code{tr\_contour} (\code{info.status}, \code{info.err}, \code{info.nodes}, \code{info.level}; the warning \code{tr\_contour:status}), and with \code{'auto'} and a fixed start in cases \Fu, \Ft\ the error estimate and status of each value returned, with the counts \code{info.ncontour}, \code{info.nfallback} and \code{info.ntail}.
\item[{[f,info] = ltratio\_pdf(x,cas,c,opts)}]
The limiting density; arguments and options as for \code{ltratio\_cdf}.  In case \Fk\ with $c_0>0$ the density is infinite at $x=-\sqrt{2c_0}$ (a logarithmic singularity, where the two zeros of the kernel's edge merge) and \code{Inf} is returned there.
\item[{[q,info] = ltratio\_icdf(p,cas,c,opts)}]
The limiting $p$-quantiles, $p\in(0,1)$, by Newton's method on \code{ltratio\_cdf} with \code{ltratio\_pdf} as derivative, safeguarded by bisection, started at the normal approximation from the first two moments.  \code{opts.tol} ($10^{-10}$) is the tolerance on $q$; \code{info.F} is the distribution function at $q$, \code{info.iter} the iteration count and \code{info.conv} the convergence flag (a failure raises \code{ltratio:icdf}).
\item[{[m,info] = ltratio\_mom(p,cas,c,opts)}]
The moments $\E(\tau^j)$, $j=1,\ldots,p$, of the limit $\tau=(U-c_0V)/\sqrt V$, as one-dimensional integrals of the joint moment generating function at real arguments (the moments of the limit in the Supplementary Material), through \code{tr\_limmom}; the options \code{c0} and \code{g0} apply, \code{method} does not.  \code{info} holds the error estimate of the integral per moment and the time.
\end{description}

\paragraph{The representations.}  With $U=\int_0^1\tilde J\dr J$ and $V=\int_0^1\tilde J^2\dr r$, $J$ the Ornstein--Uhlenbeck process $\dr J=cJ\dr r+\dr W$ started at $\gamma_0$ and $\tilde J$ its residual on the deterministic terms, the limit is $\tau=(U-c_0V)/\sqrt V$, and the joint moment generating function of $(U,V)$ is $e^{-s/2}D(s,t)^{-1/2}$ with $D$ a polynomial of degree one (\Fk) or two in $s$ (the joint moment generating function of the paper with a zero start), to whose logarithm a fixed start adds $\gamma_0^2Q(s,t)$ in closed form (the joint moment generating function with a fixed start).  The \emph{line integral representation} (\code{method = 'single'} or \code{'line'}) inverts in $s$ explicitly, which leaves a scaled $\chi^2_1$ density (\Fk) or a Bessel $K_0$ density of $U+\frac12$ under the tilt $e^{tV}$, and integrates its Gaussian transforms along the line $\Rea t=\tau$; its engine is \code{dfx\_core}.  The \emph{conditional integral representation} (\code{method = 'double'} or \code{'conditional'}) keeps the outer integral over the value $v$ of $V$ and inverts in $t$, for each $v$, along a path through a saddle point; its engine is \code{tr\_limdist}.  A fixed start $\gamma_0\ne0$ makes the tilted density noncentral and, in cases \Fu\ and \Ft, without a closed form, and these programs have no line integral there (the C and Julia implementations carry its noncentral kernel, a triple integral in cases \Fu, \Ft).  In case \Fk\ the conditional integral is then used whatever \code{method} says.  In cases \Fu\ and \Ft\ the \emph{contour integral representation} (\code{method = 'contour'}) serves, the choice of \code{'auto'} (and of \code{'single'}) there as in the C and Julia implementations: the inversion formula on paths around the singularities of the moment generating function, with the integral over $v$ in closed form, the line in $s$ replaced by two U-shaped paths around the roots of $D(\cdot,t)$ and the line in $t$ by a hyperbola, a double integral with an explicit integrand; its engine is \code{tr\_contour}, and \code{'double'} remains available.  Asked for elsewhere, \code{'contour'} is an error.  Where two representations are available they agree to about $10^{-12}$ in the distribution function (\code{ltratio\_test}, Section~\ref{sec:validation}).  The case \code{'DG'} is the GLS-detrended trend case of Elliott, Rothenberg and Stock (1996), detrended at $\bar c=-13.5$ (the DF-GLS tests of the paper): a zero start, the line integral representation only, and no moments; its numerator is of rank one, so it runs on the kernel of case \Fk\ with the Girsanov forms of case \Ft.  By default (\code{opts.verified = false}, November 2026) every value is one evaluation by one method: no second evaluation, no refinement driven by an error estimate, no fallback to another representation and no recomputation, and \code{ltratio\_method} takes the fastest representation that one evaluation gives to well below $10^{-5}$, by the timings of these programs on one thread: with a fixed start the contour integral in cases \Fu, \Ft, except where it fails ($c>1$ with $c_0<0$, its far upper tail; in case \Fu\ $c>1.5$ with $c_0>3$), and the conditional integral there and in case \Fk; with a zero or stationary start the line integral at $c_0=0$ (except in cases \Fk, \Fu\ for $-30\le c\le-5$, where the conditional integral is $1.7$ to $2.3$ times faster here), with a zero start and $c_0>0$ (where one evaluation of the conditional integral is off by up to $1.5\times10^{-2}$ at $x<0$ for $c<0$, a fault that the checked computation only flags; but the conditional integral at $x<0$ for $c\ge0$ in cases \Fu, \Ft), in case \Fk\ with $c_0<0$ for $c>-5$, in case \Sk\ for $c>-0.5$ and far from the unit root, and the conditional integral otherwise.  The C and Julia implementations differ in case \Fk\ with a fixed start, where they have the line integral and take it, at $c_0=0$, where their line integral is the faster one, and the C implementation also in cases \Su, \St, where its line integral is as fast as its conditional integral or faster.  Where the parameters lie in a regime in which one evaluation is known to be brittle the warning \code{ltratio:brittle} names it (\code{ltratio\_brittle}, decided from the parameters alone; once per call of \code{ltratio\_icdf}, \code{coef\_cdf}, \code{coef\_icdf}).  The rest of this paragraph describes \code{opts.verified = true}, the checked computation, at $1$ to $3.5$ times the cost of the default on one thread (about $1.7$ with the contour integral; more in the tails below $10^{-3}$, which it recomputes): with \code{method = 'auto'} \code{ltratio\_method} then chooses point by point the representation that is the faster one in these programs, by their own timings (\code{validation/method\_sweep.m}): the line integral at $c_0=0$, where it is $1.2$ to $60$ times faster; with $c_0\ne0$ the conditional integral in cases \Fk, \Fu, \Ft\ and \Sk, except for $c_0>0$ at $x\ge0$, where the line integral has no breakpoint and is $3$ to $30$ times faster (with $c_0<0$ its saddle factor costs up to $11$ seconds at $x$ of $4$ to $8$, while the conditional integral stays below half a second); and the line integral in cases \Su\ and \St, whose conditional integral costs half a second a value, except for $c_0<0$ at $x\ge0$.  The C and Julia implementations apply the same rule for every start; the exception for cases \Su\ and \St\ at $x<0$ is the one place where the MATLAB timings differ.  Far from the unit root, with a stationary start at $c<-50$ and with $c_0\ne0$ and a zero start at $c<-100$, all three implementations take the line integral alone, without the recomputation of the tails (Section~\ref{sec:validation}): the conditional integral is the first to lose accuracy there.  A point at which the conditional integral reports that its error estimate is not met (\code{info.status}) gets the line integral's value instead, in all three implementations (\code{info.nfallback} counts them); with \code{method = 'double'} the flagged value is returned, with a warning.  With a fixed start in cases \Fu, \Ft\ \code{'auto'} takes the contour integral at every point; a value it flags gets the conditional integral's value where that is clean, and its values below $10^{-3}$ (a density, or a distribution function in the lower tail) are recomputed by the conditional integral where that is clean, as in the C and Julia implementations, but in these programs only where the conditional integral is checked with a fixed start ($c\ge-20$, or $c\ge-40$ with $(c\gamma_0)^2/2\le100$: at $c=-30$, $\gamma_0=0.6$ in case \Fu\ it is off by $2\times10^{-9}$ relative at $F=10^{-4}$ without a flag).  Elsewhere the contour value stays, and a flagged one keeps its status in \code{info.status} (these programs have no line integral to fall back on).

\section{The engines}
\label{sec:engines}

\paragraph{The line integral: \code{dfx\_core}.}  Called through \code{dfx\_cdf(q,model,c,opts)} and the front ends \code{dfx\_pdf}, \code{dfx\_quantile}, \code{dfx\_moments}, \code{dfx\_tstar}, with the model number $1$, $2$, $3$ for \textsf{K}, \textsf{U}, \textsf{T} and \code{opts.start} equal to \code{'zero'} (default) or \code{'stationary'}; \code{ltratio\_dist} does the mapping.  The port follows the Julia implementation \code{DFExact.jl} of the line integral function by function.  The line is $\Rea t=\tau$ with $\tau=1$ for $c\le0$ and a zero start and $\min(1,t_*/2)$ otherwise, $t_*$ the convergence abscissa of the moment generating function of $V$ (\code{dfx\_tstar}; located from the cancellation-free Girsanov forms for $c>2$, so that $c$ up to $20$ and beyond is reached; for $c<-20$ it lies above $c^2/2$ and is bracketed in steps of $0.05$ in $\phi=(2t-c^2)^{1/2}$, in which the zeros of $D(0,t)$ are about $\pi$ apart).  The integrand is split, by an entire partition of unity inside the inner transforms, into a smooth part, integrated by the trapezoidal rule in $w=\tau\sinh u$ with the step $0.075$ up to $w=4900$ (for $c<-20$ the end is moved out and the step reduced, and the Filon panels are at most $2\pi|c|$ long, half a period of the factor left in $G_j$, so that the accuracy holds down to $c=-480$), and oscillating parts $G_j(w)e^{-\iu wv_j^2}$, one for each positive root $v_j$ of $X(v)=c_0v^2+xv+\frac12$ (one root, $1/(2|x|)$, at $c_0=0$ and $x<0$), integrated by Filon's method on geometric panels with $20$ Gauss--Legendre nodes; for $c>0$, and in the upper tail with $c_0<0$, the smooth part carries a saddle-point factor whose phase grows like $x^2$, and it is integrated on Gauss--Legendre panels whose node counts follow the phase.  The inner transforms are Gauss--Legendre integrals on complex paths on which the Gaussian kernel decays; with $c_0\ne0$ the paths between the breakpoints (the roots of $X$, and the vertex of the parabola for $c_0>0$) are straight pieces along the directions of steepest descent, joined through the real axis where a direct connection would cross a region in which the Gaussian grows (the line integral representation in the section on numerical evaluation in the Supplementary Material).  Options (for checking): \code{tau} (the line), \code{hu} ($0.075$), \code{N} ($20$), \code{pfac} (subdivides the Filon panels), \code{glfac} (multiplies the node counts of the Gauss--Legendre route), \code{wmax} ($4900$), \code{gen} (\code{true}: the general kernel also at $c_0=0$), \code{uniform} (the trapezoidal rule for the whole integrand); \code{dfx\_core('nodefac',2)} doubles every node count until reset with \code{1}.  The moments of \code{dfx\_moments} use the determinant form of the joint moment generating function and lose accuracy for $c>2$; \code{ltratio\_mom} uses \code{tr\_limmom}.

\paragraph{The conditional integral: \code{tr\_limdist}.}  \code{[F,info] = tr\_limdist(x,cas,c,opts,mode)}, \code{mode} \code{'cdf'} or \code{'pdf'}, with \code{opts.c0} (default $0$, as in \code{ltratio\_cdf}) and \code{opts.g0}.  In case \Fk\ it uses the first part of the conditional integral representation, the probability of a region of the plane, whose inner probability is an error function (with a fixed start, the probability $\PP[X^2\le2a+\gamma_0^2]$ of a noncentral normal); with a zero start in cases \Fu, \Ft\ and with a stationary start in all three cases the second part, the form for the $t$-ratio, whose integrands are in closed form; and in cases \Fu, \Ft\ with a fixed start the first part with the noncentral probability (\code{Pnc}), a one-dimensional integral with \code{Nl} ($96$) Gauss--Legendre nodes, so a triple integral in all.  With a zero start the integrands are the Girsanov forms (\code{iipar}), analytic beyond $t_*$, and the inner inversions run through the saddle point of the whole integrand on two paths (a parabola turning into a wedge; a vertical line in the $\vartheta$-plane; in case \Fk\ a third route, $\operatorname{erf}=1-\operatorname{erfc}$, inverts the density of $V$ and the remainder separately), with $96$ trapezoidal nodes after a sinh substitution, doubled up to three times until the nested-rule estimate is below $10^{-13}$; a value is accepted when two paths agree to that tolerance, or when the other paths are consistent with the best one within ten times the sums of the estimates (the nested rule alone being open to aliasing, \code{agreepaths}); otherwise the node counts are doubled up to three more times, after which the discrepancy is reported as the error.  With an initial observation the path is the Talbot parabola through the saddle point of $\log M_V(t)-tv$, with the sign conditions of the conditional integral representation checked at every node (\code{info.nbad} counts violations) and the branch of $\log D$ continued along it; \code{Nv} ($96$) nodes, doubled up to twice where the nested estimate exceeds $10^{-9}$ of the largest column of the piece and $10^{-11}$ divided by the range of $v$ (\code{info.nunres} counts the columns still unresolved).  The initial-value terms come from \code{tr\_nsstart} (the joint moment generating function with a fixed start: $Q(s,t)$ in closed form, from Nabeya and S\o rensen's determinants in cases \Fu, \Ft\ and the Riccati solution \code{tr\_jbstart} in case \Fk); the squared means of the noncentral form are read off $Q$ at three values of $s$.  The outer integral over $v$ is a cosine-mapped Gauss--Legendre rule in $\log v$ over $[e^{-8},e^6]\,\E V$ with \code{M} ($64$; $32$ in case \Fk\ with a zero start) nodes per part of length at most \code{plen} ($2$), breakpoints at the zeros of $a(v)=\frac12+x\sqrt v+c_0v$, grading toward the step of the integrand in the far upper tail for $c>c_0$ (width below \code{stepw}, $0.1$) and bisection, up to \code{maxlev} ($12$) times, where the Legendre-coefficient estimate exceeds \code{otol} ($10^{-12}$).  For $c<-20$ the lower end in cases \Fu, \Ft\ with a zero start is the smaller of $3\times10^{-3}$ and $e^{-3}\E V$ (deeper down the saddle-point search loses its footing, and nothing is lost above), and $|a(v)|$ is kept at least $2^{-54}$, its rounding error, away from its zeros.  \code{info.err} is the error estimate per point and \code{info.status} is $0$ when clean, otherwise $1$ (sign violations) plus $2$ (unresolved inner integrals, or \code{info.err} above \code{errtol}, $10^{-10}$); a nonzero status raises the warning \code{tr\_limdist:status}.  Further options: \code{blo}, \code{bhi} (the range in $\log v$ relative to $\log\E V$), \code{region} (checks the sign conditions on that many intermediate paths), \code{oldDK}, \code{oldII} (the earlier forms through the saddle point of $M_V$ alone, which fail for $c$ above about $2$ and $5$).

\paragraph{The contour integral: \code{tr\_contour}.}  \code{[F,info] = tr\_contour(x,cas,c,opts,mode)}, \code{cas} \code{'DU'} or \code{'DT'}, \code{mode} \code{'cdf'} or \code{'pdf'}, with \code{opts.g0} (required, nonzero) and \code{opts.c0} (default $0$); a port, step by step, of \code{tr\_contour} of \code{DFExact.jl}.  The density and the distribution function are $(2\pi\iu)^{-1}\int_C I(t)\dr t$ and $(2\pi\iu)^{-1}\int_C I(t)\dr t/t$, $I(t)=(2\pi\iu)^{-1}\{\int_{G_+}\tilde M W_+\dr s+\int_{G_-}\tilde M W_-\dr s\}$ with $\tilde M(s,t)=P(t)\{(s_1-s)(s-s_2)\}^{-1/2}\exp\{\gamma_0^2Q(s,t)\}$, $s_1$, $s_2$ the roots of $D(\cdot,t)$ and $W_\pm$ the integrals over the parts of $v>0$ where $X(v)=c_0v^2+xv+\frac12$ is positive and negative of $2v^2e^{-tv^2-sX(v)}$ (density) or $X'(v)e^{-tv^2-sX(v)}$ (distribution function), incomplete Gaussian integrals in closed form (the complex \code{erfcx} from \code{faddeeva} for arguments below $6$ in modulus, the asymptotic series beyond); the distribution function adds the kernels $(1-e^{-s/2})/s$ on $G_+$ and $-1/s$ on $G_-$.  A part whose Gaussian factor is below $e^{-50}$ from its start is dropped at that $t$.  The outer contour $C$ is the hyperbola $t(u)=\tau+a(\cosh u-1)+\iu b\sinh u$ through the line $\Rea t=\tau$ of the line integral, at the asymptotic angle of $45$ degrees for $c\ge-10$, rising to $70$ degrees at $c=-40$, integrated by the trapezoidal rule in $u$ up to $|t-\tau|=10^5$ or until a block of $32$ nodes whose terms stay below $10^{-18}$; the roots, $P(t)$, $t_*$ and $\tau$ come from \code{dfx\_core}, continued along $C$ on eight points per step.  At each node the inner paths are U-shapes around the roots, $G_+$ around $s_1$ opening to the right and $G_-$ around $s_2$ opening to the left, whose vertical segment is placed by minimizing a bound of the largest term of the quadrature plus a grading cost (a scan of $24$ distances, then $12$ steps of golden section), graded by Gauss--Legendre panels toward the origin, the other root and the branch point $-t/c_0$, and continued by horizontal rays.  By default (\code{opts.verified = false}) the value is evaluation A alone at its first level ($64$ nodes; \code{info.err} \code{NaN}, \code{info.status} $0$).  With \code{opts.verified = true} the error estimate \code{info.err} comes from two structurally different evaluations, A as described with $64$ nodes per panel and ray, and B with the asymptotic angle of the hyperbola larger by $7.5$ degrees (so that every $t$ differs), the grading cost weighted by $1$ instead of $\frac12$ in the placement (so that the U-shapes and their panels differ) and $48$ nodes: it is $|A-B|$ plus $10\epsilon$ times the sums of the largest terms of both, and while it exceeds $10^{-10}$ the next level (A: $96$, $128$, then $192$ nodes; B: $64$, $96$, $128$) is added to both at all their nodes; the value is A at the last level (\code{info.level}, $1$ to $4$).  A value is thus accepted only when two evaluations on different paths with different rules agree, which a path in the wrong place, whose error a finer rule on the same path would reproduce, does not pass.  \code{info.status} is $0$ when the estimate is at most $10^{-10}$, otherwise $2$, with the warning \code{tr\_contour:status}, and \code{info.nodes} is the number of outer nodes of the two evaluations.

\paragraph{The moment generating function.}  \code{tr\_cflim(t1,t2,cas,c,form)} evaluates the Nabeya--Tanaka determinants of the joint moment generating function with a zero start at complex arguments (31-term Taylor series in $z$ for $|z|<4$ from \code{nt\_series}, closed forms in $\vartheta=(-z)^{1/2}$ and $e^{-2\vartheta}$ beyond), with the branch of $\log D$ followed along two legs and the roots of $D$ in $s$ (the weights $\kappa_k$ of $U+\frac12$ under the tilt); the forms \code{'F'}, \code{'B'}, \code{'BC'} and \code{'ApB'} return the coefficients, the weights, the branch-continued logarithm along a path, and the moment generating function of $U$ and $V$ together.  \code{tr\_limmom(p,cas,c,opts)} computes the moments by Cauchy integrals on a circle in the argument of $U$ and the trapezoidal rule in $\log s$, from the Girsanov forms with a zero start and with the initial-value terms of \code{tr\_nsstart} otherwise.

\section{Accuracy and cost}
\label{sec:accuracy}

Both representations are accurate to about $10^{-11}$ in the distribution function and $10^{-10}$ relative in the density over $c\in[-30,20]$, $c_0\in[-40,10]$ and the six cases, and to $10^{-11}$ with a fixed start in cases \Fu, \Ft\ for $c\le2$ (the section on numerical evaluation in the Supplementary Material).  With a zero or stationary start the line integral keeps this accuracy down to $c=-480$, where the conditional integral does not: at $c=-60$, $-100$, $-200$ and $-480$ these programs agree with the Julia and C implementations to $8\times10^{-12}$, and \code{ltratio\_test} checks them against the Julia reference values at $c=-100$ and $-480$.  With a fixed start $\gamma_0\ne0$ the conditional integral is checked for $c\ge-20$ only (case \Fk, and \code{'double'} in cases \Fu, \Ft); the contour integral covers the region of the C and Julia implementations in cases \Fu\ and \Ft\ ($-100\le c\le2$, $|\gamma_0|\le1.5$ and $(c\gamma_0)^2/2\le4050$ in case \Fu, $|\gamma_0|\le3$ and $(c\gamma_0)^2/2\le7200$ in case \Ft), and \code{ltratio\_test} checks it against thirty reference values of the Julia implementation, from $c=-94.4$ to $2$, to $10^{-11}$.  The error of the line integral is absolute rather than relative: its terms are of order one, and it carries an error of $10^{-14}$ to $10^{-21}$ whatever the size of the value, which finer rules do not reduce ($2\times10^{-14}$ in the lower tails of cases \Ft\ and \St\ with $c_0\ne0$).  With \code{opts.verified = true} and \code{opts.method} not given, \code{ltratio\_pdf} and \code{ltratio\_cdf} therefore recompute by the conditional integral every value of the line integral below $10^{-3}$ (a density, or a distribution function in the lower tail), with its error estimate in \code{info.err}, except far from the unit root, where the line integral serves alone (a stationary start at $c<-50$; $c_0\ne0$ with a zero start at $c<-100$).  The relative accuracy then holds in the tails as well, down to that estimate; far out, below the error of either method, a computed density can be a tiny negative number, and it is returned as zero.  Doubling every node count of the line integral changes the distribution function by at most $4\times10^{-13}$; the two representations agree to $4\times10^{-12}$ at worst (\code{validation/acc\_table.m}).  With a fixed start in cases \Fu, \Ft\ at $c$ above about $2.5$ the Talbot path of the conditional integral no longer satisfies the sign conditions everywhere, and \code{tr\_limdist} warns; the line integral does not cover that configuration.  A value of the distribution function or the density takes, in MATLAB on one thread of a laptop (the default, one evaluation), $0.03$ to $0.1$ seconds by the line integral for $c_0=0$ and $0.2$ to $1.1$ seconds with $c_0\ne0$ (the upper end in the upper tail at $c=5$, $c_0=-5$), $0.02$ to $0.11$ seconds by the conditional integral with a zero start, $0.4$ to $0.5$ seconds with a stationary start and $1.3$ to $2.1$ seconds with a fixed start in cases \Fu, \Ft, and $0.05$ to $0.14$ seconds by the contour integral with a fixed start in cases \Fu, \Ft; a quantile takes a few values and a set of four moments about a second.  The Julia implementation, \code{DFExact.jl}, takes $0.5$ to $10$ milliseconds per value by the line integral at $c_0=0$ and $7$ to $15$ milliseconds by the contour integral, on ten threads.

\section{Validation and reproduction of the paper}
\label{sec:validation}

\code{ltratio\_test} runs two hundred and sixty-two checks in about five minutes and reports a failure count: the two representations against each other at the mean plus $-2.5$ to $2.5$ standard deviations in the six cases, at $c\in\{0,-5,-10\}$, $c=2$ and $5$ and $c_0\ne0$; the exact values $\PP[\hat\tau\le0]=\operatorname{erf}\{(2w)^{-1/2}\}$, $w=(e^{2c}-1)/(2c)$, for $c\in\{-10,0,2,5,20\}$, $\PP[\hat\tau_\mu\le0]=0.95756559$ from the coefficient statistic, and $\PP[U-c_0V\le0]$ at $c=10$, $c_0=5$ from an independent one-dimensional inversion; the distribution function at the asymptotic critical values of \citet{MacKinnon2010} and the quantile function; the initial observation (continuity as $\gamma_0\to0$; the stationary start as the 24-point Gauss--Hermite average of the fixed start; the noncentral probability $\PP[J(1)^2\le1+\gamma_0^2]$ in case \Fk; the rank-two structure of the initial-value term); the moments against \citet{Nabeya1999}, against \code{dfx\_moments} and against the integral of the density; the density against fifth-order differences of the distribution function; the distribution function, the density, the moments and the quantiles far from the unit root, at $c=-100$ and $-480$, against the reference values of the Julia implementation; the contour integral against thirty reference values of the Julia implementation (a fixed start in cases \Fu, \Ft), the automatic choice there and its scope; and the argument checks.  \code{validation/acc\_table.m} produces the accuracy table and its LaTeX rows.

The reproduction of the paper is not part of this package: it is in the paper's repository, in the folder \code{public/replication}, and runs on these programs.  There \code{ltratio\_paper(which,opts)} recomputes the tables of the paper, the right-tail $p$-values at the unit root (with \code{mackinnon94\_p}, the approximate $p$-values of MacKinnon (1994) that statsmodels and Stata report), the numbers quoted on local power, confidence sets and the weighted-average-power test, and the data of the figures, and prints each value next to the published one; the parts are \code{'quant'} (alias \code{'adfmax'}), \code{'larsson'}, \code{'init'}, \code{'mom'}, \code{'bubble'}, \code{'pvalues'}, \code{'power'}, \code{'ci'} and \code{'wap'}.  The data of the figures are compared with the files in \code{LaTeX/figures} of that repository at every \code{opts.step}-th point (default $5$), and with \code{opts.write = true} computed on the full grid and written anew in the same format.  It needs the folder \code{matlab} of this package on the path and stops with a message when \code{ltratio\_cdf} is not found.  The Julia program \code{ltratio\_paper.jl} in \code{julia/replication} of the paper's repository has the same parts, options and output.

\section{Files}
\label{sec:files}

\begin{description}
\item[ltratio\_cdf.m, ltratio\_pdf.m, ltratio\_icdf.m, ltratio\_mom.m] the user-level routines; \code{ltratio\_dist.m} (dispatch to the engines), \code{ltratio\_method.m} (the automatic choice of the representation) and \code{ltratio\_args.m} (argument checks).
\item[dfx\_core.m, dfx\_cdf.m, dfx\_pdf.m, dfx\_quantile.m] the line integral: the engine and its front ends;
\item[dfx\_moments.m, dfx\_tstar.m, dfx\_series.m, dfx\_test.m] its moments, convergence abscissa, series coefficients and test.
\item[tr\_limdist.m, tr\_limmom.m] the conditional integral and the moments;
\item[tr\_contour.m] the contour integral (a fixed start in cases \Fu, \Ft);
\item[ltratio\_brittle.m] the warning of the one-evaluation default in a regime known to be brittle;
\item[tr\_cflim.m, tr\_nsstart.m, tr\_jbstart.m, nt\_series.m, gen\_series.py, faddeeva.m] the moment generating function, the initial-value terms, the series coefficients and the Faddeeva function.
\item[adfmax\_cdf.m, adfmax\_pdf.m, adfmax\_icdf.m, coef\_cdf.m, coef\_icdf.m] the ADFmax and the coefficient statistics of the paper;
\item[ltratio\_region.m, ltratio\_lincdf.m, power\_envelope.m] regions of the plane, linear forms and the power envelope (the applications of the paper), with the engines \code{tr\_adfmax.m}, \code{tr\_region.m}, \code{tr\_gilpelaez.m} and \code{ltratio\_bisect.m}, \code{gls\_parts.m}, \code{lam\_gls.m}.
\item[ltratio\_test.m, validation/] validation.
\item[Contents.m, README.md, LICENSE, doc/] this manual.
\end{description}

\begin{thebibliography}{9}
\bibitem[Bao et~al.(2026a)]{BaoBrodaKan2026l}
Bao, Y., Broda, S.~A. and Kan, R. (2026a).  The exact limiting distribution of the Dickey--Fuller $t$-ratio and some related statistics.  Manuscript.
\bibitem[Bao et~al.(2026b)]{BaoBrodaKan2026}
Bao, Y., Broda, S.~A. and Kan, R. (2026b).  The exact finite-sample distribution of the Dickey--Fuller $t$-ratio.  Manuscript.
\bibitem[MacKinnon(2010)]{MacKinnon2010}
MacKinnon, J.~G. (2010).  Critical values for cointegration tests.  Queen's Economics Department Working Paper 1227.
\bibitem[Nabeya(1999)]{Nabeya1999}
Nabeya, S. (1999).  Asymptotic moments of some unit root test statistics in the null case.  \emph{Econometric Theory} 15, 139--149.
\end{thebibliography}

\end{document}
